\documentclass[11pt]{article}
\usepackage[margin=1in]{geometry}
\usepackage[T1]{fontenc}
\usepackage[utf8]{inputenc}
\usepackage{amsmath,amssymb,amsthm}
\usepackage{enumitem}

\title{ICPC World Finals 2022\\Z. Archaeological Recovery}
\author{}
\date{}

\begin{document}
\maketitle

\section*{Problem Summary}

Encode each pyramid configuration as a vector in $\mathbb{Z}_3^k$. Then each lever is also a vector
$u_j \in \mathbb{Z}_3^k$, and pulling a set of levers adds their vectors. The input gives, for each final
configuration, how many subsets of the unknown levers produce it. We must reconstruct one valid
multiset of $n$ lever vectors.

\section*{Key Observations}

\begin{itemize}[leftmargin=*]
    \item For any fixed $x \in \mathbb{Z}_3^k$, taking dot products with $x$ turns the whole problem
    into a one-dimensional problem over $\mathbb{Z}_3$.
    \item In dimension $1$, the number of zero levers is exactly the exponent of the largest power of
    $2$ dividing all three residue counts. This lets us recover
    \[
    f(x) = |\{j : u_j \cdot x = 0\}|
    \]
    for every $x$.
    \item Averaging $f$ over an orthogonal complement tells us how many levers lie in a given linear
    subspace. In particular, for each line $\{0,x,-x\}$ we can recover how many levers belong to the
    pair $\{x,-x\}$.
    \item After that, only the signs are unknown: for every pair $\{x,-x\}$ we know how many copies
    exist, but not which copies are $x$ and which are $-x$.
    \item Flipping the sign of a subset of levers translates the entire subset-sum multiset by one vector
    of $\mathbb{Z}_3^k$. So after choosing arbitrary signs, the real answer differs only by a shift of
    the generated multiset.
\end{itemize}

\section*{Algorithm}

\begin{enumerate}[leftmargin=*]
    \item Enumerate all $3^k$ vectors of $\mathbb{Z}_3^k$ and precompute:
    \begin{itemize}[leftmargin=*]
        \item vector addition modulo $3$,
        \item vector negation,
        \item dot products modulo $3$.
    \end{itemize}
    \item For each $x \in \mathbb{Z}_3^k$, project the observed subset-sum multiset onto the three
    residues of $s \cdot x$. Let those counts be $c_0,c_1,c_2$. From the one-dimensional fact above,
    set
    \[
    f(x) = \nu_2(\gcd(c_0,c_1,c_2)),
    \]
    where $\nu_2$ is the exponent of the largest power of $2$ dividing its argument.
    \item Recover the multiplicity of the zero vector by averaging $f$ over all $x$:
    \[
    m_{\{0\}} =
    \frac{1}{2}\left(\frac{3}{3^k}\sum_x f(x) - n\right).
    \]
    \item For each nonzero line $G = \{0,x,-x\}$, let $H = G^\perp$. Then recover the number of
    levers in that line by
    \[
    m_G =
    \frac{1}{2}\left(\frac{3}{|H|}\sum_{y \in H} f(y) - n\right).
    \]
    Subtract the zero-vector multiplicity to obtain the number of copies of $\{x,-x\}$.
    \item Choose one canonical representative from each pair $\{x,-x\}$ and create that many copies
    of it. Together with the zero vectors this gives an arbitrary-sign lever multiset.
    \item Compute the subset-sum multiset of this arbitrary-sign multiset by DP over all $3^k$ states.
    At the same time, store one witness subset for every reachable sum.
    \item Try every shift vector $v \in \mathbb{Z}_3^k$. If shifting the DP multiset by $v$ matches the
    target multiset, and $v$ is reachable as a subset sum of the arbitrary-sign levers, flip the signs of
    exactly the witness subset for $v$ and output the resulting levers.
\end{enumerate}

\section*{Correctness Proof}

We prove that the algorithm returns the correct answer.

\paragraph{Lemma 1.}
For every $x \in \mathbb{Z}_3^k$, the value $f(x) = |\{j : u_j \cdot x = 0\}|$ is determined by the
projected subset-sum counts onto the residues of $s \cdot x$.

\paragraph{Proof.}
Fix $x$ and replace every lever vector $u_j$ by the scalar $u_j \cdot x \in \mathbb{Z}_3$. The observed
subset sums project to the subset sums of these scalars, so we obtain a one-dimensional instance.

In that one-dimensional instance, every zero lever doubles all three counts, so each zero lever
contributes one factor of $2$ to their common gcd. If we remove all zero levers and only keep
nonzero scalars, the three counts are never all even; equivalently, the common gcd is odd. Hence
the exact number of zero scalars is the largest power of $2$ dividing all three counts, which is
precisely $\nu_2(\gcd(c_0,c_1,c_2))$. \qed

\paragraph{Lemma 2.}
For any linear subspace $G \subseteq \mathbb{Z}_3^k$, the values $f(x)$ determine how many levers
lie in $G$.

\paragraph{Proof.}
Let $H = G^\perp$. Consider
\[
\sum_{y \in H} f(y).
\]
Each lever inside $G$ is orthogonal to every $y \in H$, so it contributes $|H|$ to this sum. A lever
outside $G$ is orthogonal to exactly one third of the vectors in $H$, so it contributes $|H|/3$.

If $m_G$ levers lie in $G$, then
\[
\sum_{y \in H} f(y) = m_G |H| + (n-m_G)\frac{|H|}{3}.
\]
Solving for $m_G$ gives
\[
m_G =
\frac{1}{2}\left(\frac{3}{|H|}\sum_{y \in H} f(y) - n\right),
\]
which is exactly the formula used by the algorithm. \qed

\paragraph{Lemma 3.}
Once the multiplicity of each pair $\{x,-x\}$ is known, any valid lever multiset differs from an
arbitrary sign choice only by flipping the signs of some subset of the chosen levers. Such a sign flip
translates the whole subset-sum multiset by a reachable vector.

\paragraph{Proof.}
Knowing the multiplicity of each pair $\{x,-x\}$ means that only the sign of each nonzero copy is
still undetermined. Therefore any valid solution can be obtained from the arbitrary-sign choice by
negating some subset of those copies.

Negating one lever $u$ changes every subset sum by the same translation: the multiset generated by
$-u$ is the old multiset shifted by $-u$. Repeating this for several levers shows that negating a
subset translates the entire multiset by the sum of the originally chosen vectors in that subset.
That translation vector is reachable as a subset sum, and the stored witness subset reconstructs
exactly which signs to flip. \qed

\paragraph{Theorem.}
The algorithm outputs a valid multiset of lever vectors.

\paragraph{Proof.}
By Lemma 1, the algorithm computes the exact values $f(x)$ for all $x$. By Lemma 2, these values
recover the multiplicity of the zero vector and of every pair $\{x,-x\}$, so the only remaining
freedom is the choice of signs.

By Lemma 3, every valid sign choice corresponds to a translation of the subset-sum multiset of the
arbitrary-sign construction by some reachable vector. The algorithm tests all shifts and, once it finds
the matching one, flips exactly the witness subset that realizes that shift. The resulting lever multiset
therefore generates exactly the target subset-sum multiset. \qed

\section*{Complexity Analysis}

Let $M = 3^k$. Precomputing addition and dot products takes $O(kM^2)$ time and $O(M^2)$ memory.
Computing all values $f(x)$ and all line multiplicities takes another $O(M^2)$ time. The subset-sum
DP over the reconstructed levers takes $O(nM)$ time and $O(M)$ extra memory. Thus the total
running time is $O(kM^2 + nM)$, which is $O(3^{2k} + n3^k)$ up to constant factors, and the memory
usage is $O(M^2)$.

\section*{Implementation Notes}

\begin{itemize}[leftmargin=*]
    \item The states are encoded as base-$3$ integers, but vector addition must still be done
    coordinate-wise modulo $3$; it is not the same as ordinary integer addition modulo $3^k$.
    \item The DP stores one witness subset for every reachable sum so that the final sign flips can be
    reconstructed directly.
\end{itemize}

\end{document}
